% Answers to Chapter 15, from lectures/L15/appendix/c_solutions.md; the self-test from
% lectures/L15/README.md.

\facitavsnitt{c15}{Kapitel 15}{Komplexa tal, del I}
\begin{facit}

\svar{1.1} \sv{a} $u(t) = (3t^2 - 4t + 2)\,\text{V}$ \sv{b} $17\,\text{V}$

\svar{2.1} $U = \sqrt{29}\angle{-0{,}38}\,\text{rad} \approx 5{,}39\angle{-0{,}38}\,\text{rad}$~V,
det vill säga $\delta_u \approx -21{,}8^{\circ}$.

\svar{2.2} $I = \frac{5\sqrt{3}}{2} + j2{,}5 \approx (4{,}33 + j2{,}5)\,\text{mA}$

\svar{2.3} $Z \approx 1{,}08\angle{-0{,}9}\,\text{rad}$~k$\Omega$
$\approx (0{,}67 - j0{,}85)\,\text{k}\Omega$

\svar{3.1} \sv{a} $-1$ \sv{b} $-j$ \sv{c} $1$ \sv{d} $j$ \sv{e} $-j$ \sv{f} $x = \pm j5$

\svar{3.2} \sv{a} Realdel och imaginärdel: $z_1$: $4$ och $-3$;\enspace $z_2$: $-2$ och
$5$;\enspace $z_3$: $0$ och $6$;\enspace $z_4$: $-7$ och $0$.
\sv{b} $z_3 = j6$ \sv{c} $z_4 = -7$

\svar{3.3} \sv{a} $2 + j7$ \sv{b} $4 - j3$ \sv{c} $3 + j19$ \sv{d} $-7 + j$

\svar{3.4} \sv{a} $-10 + j11$ \sv{b} $13$ \sv{c} $2 + j5$ \sv{d} $j2$

\svar{3.5} \sv{a} $5 - j2$ \sv{b} $-3 + j7$ \sv{c} $25$
\sv{d} $(x + jy)(x - jy) = x^2 - j^2y^2 = x^2 + y^2 = \abs{z}^2$

\svar{3.6} \sv{a} $1 + j3$ \sv{b} $8 - j6$ \sv{c} $1{,}2 - j1{,}6$
\sv{d} Produkten av nämnaren och dess konjugat är reell, och med en reell nämnare kan real- och
imaginärdelen divideras var för sig.

\svar{3.7} \sv{a} $5$ \sv{b} $13$ \sv{c} $7$ \sv{d} $8$ \sv{e} $\sqrt{2} \approx 1{,}41$

\svar{3.8} \sv{a} $45^{\circ}$, kvadrant I \sv{b} $135^{\circ}$, kvadrant II
\sv{c} $225^{\circ}$ (eller $-135^{\circ}$), kvadrant III \sv{d} $-45^{\circ}$, kvadrant IV
\sv{e} $53{,}1^{\circ}$, kvadrant I

\svar{3.9} \sv{a} $10\angle 53{,}1^{\circ}$ \sv{b} $5\angle 126{,}9^{\circ}$
\sv{c} $\sqrt{2}\angle 225^{\circ} \approx 1{,}41\angle 225^{\circ}$ \sv{d} $5\angle{-90^{\circ}}$

\svar{3.10} \sv{a} $8{,}66 + j5$ \sv{b} $j4$ \sv{c} $-6$ \sv{d} $1 - j1{,}73$
\sv{e} $5{,}66 + j5{,}66$

\svar{3.11} \sv{a} $13\,\text{V}$ \sv{b} $22{,}6^{\circ}$ \sv{c} $13\angle 22{,}6^{\circ}\,\text{V}$
\sv{d} Hur mycket spänningen är förskjuten i tid relativt referensen, till exempel strömmen.

\svar{3.12} \sv{a} $Z = (30 + j40)\,\Omega$ \sv{b} $50\,\Omega$ \sv{c} $53{,}1^{\circ}$
\sv{d} Induktiv: imaginärdelen är positiv.

\svar{3.13} \sv{a} $\abs{U} = 10\,\text{V}$, $\abs{Z} = 5\,\Omega$
\sv{b} $I = (1{,}2 - j1{,}6)\,\text{A}$ \sv{c} $\abs{I} = \frac{10}{5} = 2\,\text{A}$
\sv{d} $-53{,}1^{\circ}$: strömmen ligger $53{,}1^{\circ}$ efter spänningen, som i en induktiv
krets.

\svar{3.14} \sv{a} $x = \pm j3$ \sv{b} $x = -2 \pm j3$ \sv{c} $x = 1 \pm j2$
\sv{d} De är varandras konjugat: samma realdel och motsatta imaginärdelar.

\svar{3.15} \sv{a} $8\,\Omega$. Imaginärdelarna tar ut varandra, så seriekopplingen är rent
resistiv. \sv{b} Båda är $5\,\Omega$. \sv{c} $25\,\Omega^2$ \sv{d} $3{,}125\,\Omega$

\svar{3.16} \sv{a} Falskt: $j^2 = -1$.
\sv{b} Falskt: $\abs{z}$ är en längd, alltid $\geq 0$.
\sv{c} Sant: talet ligger på den imaginära axeln.
\sv{d} Sant: produkten är $a^2 + b^2$.
\sv{e} Falskt: $\arctan$ ger bara vinklar mellan $-90^{\circ}$ och $90^{\circ}$; i kvadrant II
och III läggs $180^{\circ}$ till.
\sv{f} Falskt: addition är enklast på rektangulär form.

\svar[Testa dig själv]{T} \sv{1} $\abs{z} = 5$ och $\delta \approx 0{,}927\,\text{rad}$, alltså
$z = 5\angle 0{,}927$. \sv{2} $5 + j$

\end{facit}
